\documentclass[otchet]{SCWorks}
% Тип обучения (одно из значений):
%    bachelor   - бакалавриат (по умолчанию)
%    spec       - специальность
%    master     - магистратура
% Форма обучения (одно из значений):
%    och        - очное (по умолчанию)
%    zaoch      - заочное
% Тип работы (одно из значений):
%    coursework - курсовая работа (по умолчанию)
%    referat    - реферат
%  * otchet     - универсальный отчет
%  * nirjournal - журнал НИР
%  * digital    - итоговая работа для цифровой кафедры
%    diploma    - дипломная работа
%    pract      - отчет о научно-исследовательской работе
%    autoref    - автореферат выпускной работы
%    assignment - задание на выпускную квалификационную работу
%    review     - отзыв руководителя
%    critique   - рецензия на выпускную работу

% * Добавлены вручную. За вопросами к @mchernigin
\usepackage{preamble}

\begin{document}

% Кафедра (в родительном падеже)
\chair{математической кибернетики и компьютерных наук}

% Тема работы
\title{Анализ красно-чёрного дерева}

% Курс
\course{2}

% Группа
\group{211}

% Факультет (в родительном падеже) (по умолчанию "факультета КНиИТ")
% \department{факультета КНиИТ}

% Специальность/направление код - наименование
 \napravlenie{02.03.02 "--- Фундаментальная информатика и информационные технологии}
% \napravlenie{02.03.01 "--- Математическое обеспечение и администрирование информационных систем}
% \napravlenie{09.03.01 "--- Информатика и вычислительная техника}
%\napravlenie{09.03.04 "--- Программная инженерия}
% \napravlenie{10.05.01 "--- Компьютерная безопасность}

% Для студентки. Для работы студента следующая команда не нужна.
% \studenttitle{студентки}

% Фамилия, имя, отчество в родительном падеже
\author{Бурова Арсения Александровича}

% Заведующий кафедрой 
\chtitle{доцент, к.\,ф.-м.\,н.}
\chname{С.\,В.\,Миронов}

% Руководитель ДПП ПП для цифровой кафедры (перекрывает заведующего кафедры)
% \chpretitle{
%     заведующий кафедрой математических основ информатики и олимпиадного\\
%     программирования на базе МАОУ <<Ф"=Т лицей №1>>
% }
% \chtitle{г. Саратов, к.\,ф.-м.\,н., доцент}
% \chname{Кондратова\, Ю.\,Н.}

% Научный руководитель (для реферата преподаватель проверяющий работу)
\satitle{доцент, к.\,ф.-м.\,н.} %должность, степень, звание
\saname{Ю.\,Н.\,Кондратова}

% Руководитель практики от организации (руководитель для цифровой кафедры)
\patitle{доцент, к.\,ф.-м.\,н.}
\paname{С.\,В.\,Миронов}

% Руководитель НИР
\nirtitle{доцент, к.\,п.\,н.} % степень, звание
\nirname{В.\,А.\,Векслер}

% Семестр (только для практики, для остальных типов работ не используется)
\term{2}

% Наименование практики (только для практики, для остальных типов работ не
% используется)
\practtype{учебная}

% Продолжительность практики (количество недель) (только для практики, для
% остальных типов работ не используется)
\duration{2}

% Даты начала и окончания практики (только для практики, для остальных типов
% работ не используется)
\practStart{01.07.2024}
\practFinish{13.01.2024}

% Год выполнения отчета
\date{2026}

\maketitle

% Включение нумерации рисунков, формул и таблиц по разделам (по умолчанию -
% нумерация сквозная) (допускается оба вида нумерации)
\secNumbering

\tableofcontents

% Раздел "Обозначения и сокращения". Может отсутствовать в работе
% \abbreviations
% \begin{description}
%     \item ... "--- ...
%     \item ... "--- ...
% \end{description}

% Раздел "Определения". Может отсутствовать в работе
% \definitions

% Раздел "Определения, обозначения и сокращения". Может отсутствовать в работе.
% Если присутствует, то заменяет собой разделы "Обозначения и сокращения" и
% "Определения"
% \defabbr

\section{Реализация}

\begin{Verbatim}[
  numbers=left,
  breaklines=true,
  breakanywhere=true
]

tree* node(tree* p, int x) {
    tree* n = new tree;
    n->inf = x;
    n->color = RED;
    n->parent = p;
    n->left = n->right = NULL;
    return n;
}

tree* root(int x) {
    tree* n = new tree;
    n->inf = x;
    n->color = BLACK;
    n->parent = n->left = n->right = NULL;
    return n;
}

//вспомогательные функции для работы с деревом
tree* grandparent(tree* n) {
    if (n && n->parent && n->parent->parent)
        return n->parent->parent;
    return NULL;
}

tree* uncle(tree* n) {
    tree* g = grandparent(n);
    if (!g) return NULL;
    if (n->parent == g->left)
        return g->right;
    else
        return g->left;
}

tree* sibling(tree* n) {
    if (!n || !n->parent) return NULL;
    if (n == n->parent->left)
        return n->parent->right;
    else
        return n->parent->left;
}

//поворот влево
void rotate_left(tree*& tr, tree* n) {
    if (!n || !n->right) return;

    tree* y = n->right;
    n->right = y->left;

    if (y->left)
        y->left->parent = n;

    y->parent = n->parent;

    if (!n->parent)
        tr = y;
    else if (n == n->parent->left)
        n->parent->left = y;
    else
        n->parent->right = y;

    y->left = n;
    n->parent = y;
}

//поворот вправо
void rotate_right(tree*& tr, tree* n) {
    if (!n || !n->left) return;
    tree* y = n->left;
    n->left = y->right;

    if (y->right)
        y->right->parent = n;

    y->parent = n->parent;

    if (!n->parent)
        tr = y;
    else if (n == n->parent->right)
        n->parent->right = y;
    else
        n->parent->left = y;

    y->right = n;
    n->parent = y;
}

//функции на разные случаи вставки
void insert_case5(tree*& tr, tree* n) {
    tree* g = grandparent(n);
    if (!g) return;

    n->parent->color = BLACK;
    g->color = RED;

    if (n == n->parent->left) {
        rotate_right(tr, g);
    }
    else {
        rotate_left(tr, g);
    }
}

void insert_case4(tree*& tr, tree* n) {
    tree* g = grandparent(n);
    if (!g) return;

    if (n == n->parent->right && n->parent == g->left) {
        rotate_left(tr, n->parent);
        n = n->left;
    }
    else if (n == n->parent->left && n->parent == g->right) {
        rotate_right(tr, n->parent);
        n = n->right;
    }

    insert_case5(tr, n);
}

void insert_case3(tree*& tr, tree* n) {
    tree* u = uncle(n);

    if (u && u->color == RED) {
        n->parent->color = BLACK;
        u->color = BLACK;
        tree* g = grandparent(n);
        if (g) {
            g->color = RED;
            insert_case1(tr, g);
        }
    }
    else {
        insert_case4(tr, n);
    }
}

void insert_case2(tree*& tr, tree* n) {
    if (!n->parent) return;
    if (n->parent->color == BLACK)
        return;
    else
        insert_case3(tr, n);
}

void insert_case1(tree*& tr, tree* n) {
    if (!n->parent) {
        n->color = BLACK;
        tr = n;
    }
    else {
        insert_case2(tr, n);
    }
}

//функция для вставки
void insert(tree*& tr, int x) {
    if (!tr) {
        tr = root(x);
        return;
    }

    tree* y = tr;
    tree* p = nullptr;

    while (y) {
        p = y;
        if (x < y->inf)
            y = y->left;
        else if (x > y->inf)
            y = y->right;
        else
            return;  //элемент уже существует
    }

    tree* n = node(p, x);
    if (x < p->inf)
        p->left = n;
    else
        p->right = n;

    insert_case1(tr, n);
}

//функция для создания дерева
tree* create_tree() {
    tree* tr = NULL;
    int n, x;
    cout << "введите количество элементов: ";
    cin >> n;

    if (n > 0) {
        cout << "элемент 1: ";
        cin >> x;
        tr = root(x);

        for (int i = 1; i < n; i++) {
            cout << "элемент " << i + 1 << ": ";
            cin >> x;
            insert(tr, x);
        }
    }
    return tr;
}

//функции для удаления
void replace_node(tree*& tr, tree* old_node, tree* new_node) {
    if (!old_node->parent) {
        tr = new_node;
        if (new_node) new_node->parent = NULL;
    }
    else {
        if (old_node == old_node->parent->left)
            old_node->parent->left = new_node;
        else
            old_node->parent->right = new_node;

        if (new_node)
            new_node->parent = old_node->parent;
    }
}

void delete_case6(tree*& tr, tree* n) {
    if (!n || !n->parent) return;

    tree* s = sibling(n);
    if (!s) return;

    s->color = n->parent->color;
    n->parent->color = BLACK;

    if (n == n->parent->left) {
        if (s->right) s->right->color = BLACK;
        rotate_left(tr, n->parent);
    }
    else {
        if (s->left) s->left->color = BLACK;
        rotate_right(tr, n->parent);
    }
}

void delete_case5(tree*& tr, tree* n) {
    if (!n || !n->parent) {
        delete_case6(tr, n);
        return;
    }

    tree* s = sibling(n);
    if (!s) {
        delete_case6(tr, n);
        return;
    }

    if (s->color == BLACK) {
        if ((n == n->parent->left) &&
            (!s->right || s->right->color == BLACK) &&
            (s->left && s->left->color == RED)) {

            s->color = RED;
            s->left->color = BLACK;
            rotate_right(tr, s);
        }
        else if ((n == n->parent->right) &&
            (!s->left || s->left->color == BLACK) &&
            (s->right && s->right->color == RED)) {

            s->color = RED;
            s->right->color = BLACK;
            rotate_left(tr, s);
        }
    }

    delete_case6(tr, n);
}

void delete_case4(tree*& tr, tree* n) {
    if (!n || !n->parent) {
        delete_case5(tr, n);
        return;
    }

    tree* s = sibling(n);
    if (!s) {
        delete_case5(tr, n);
        return;
    }

    if (n->parent->color == RED &&
        s->color == BLACK &&
        (!s->left || s->left->color == BLACK) &&
        (!s->right || s->right->color == BLACK)) {

        s->color = RED;
        n->parent->color = BLACK;
    }
    else {
        delete_case5(tr, n);
    }
}

void delete_case3(tree*& tr, tree* n) {
    if (!n) return;

    tree* s = sibling(n);
    bool parentBlack = (!n->parent || n->parent->color == BLACK);
    bool siblingBlack = (!s || s->color == BLACK);
    bool leftNephewBlack = (!s || !s->left || s->left->color == BLACK);
    bool rightNephewBlack = (!s || !s->right || s->right->color == BLACK);

    if (parentBlack && siblingBlack && leftNephewBlack && rightNephewBlack) {
        if (s) s->color = RED;
        if (n->parent)
            delete_case1(tr, n->parent);
    }
    else {
        delete_case4(tr, n);
    }
}

void delete_case2(tree*& tr, tree* n) {
    if (!n || !n->parent) {
        delete_case3(tr, n);
        return;
    }

    tree* s = sibling(n);

    if (s && s->color == RED) {
        n->parent->color = RED;
        s->color = BLACK;

        if (n == n->parent->left)
            rotate_left(tr, n->parent);
        else
            rotate_right(tr, n->parent);
    }

    delete_case3(tr, n);
}

void delete_case1(tree*& tr, tree* n) {
    if (!n) return;

    if (!n->parent) {
        if (n->left) {
            tr = n->left;
            tr->color = BLACK;
            tr->parent = NULL;
        }
        else if (n->right) {
            tr = n->right;
            tr->color = BLACK;
            tr->parent = NULL;
        }
    }
    else {
        delete_case2(tr, n);
    }
}

tree* find_min(tree* node) {
    if (!node) return NULL;
    while (node->left)
        node = node->left;
    return node;
}

tree* find_max(tree* node) {
    if (!node) return NULL;
    while (node->right)
        node = node->right;
    return node;
}

//основная функция удаления
void delete_node(tree*& tr, int x) {
    if (!tr) return;  //дерево пустое

    //поиск удаляемого узла
    tree* n = tr;
    while (n) {
        if (x < n->inf)
            n = n->left;
        else if (x > n->inf)
            n = n->right;
        else
            break;
    }
    if (!n) {
        cout << "Узел не найден!" << endl;
        return;
    }

    //случай 1: удаляемый узел имеет двух реальных детей
    if (n->left && n->right) {
        tree* buf;

        //выбор замены
        if (n->inf >= tr->inf) {
            buf = find_max(n->left);
        }
        else {
            buf = find_min(n->right);
        }

        if (!buf) return;

        n->inf = buf->inf;
        n = buf;
    }

    //определяем ребенка
    tree* child = n->left ? n->left : n->right;

    //случай 2: удаляемый узел имеет одного ребенка
    if (child) {
        //заменяем узел его ребенком
        replace_node(tr, n, child);

        //если удаляемый узел черный
        if (n->color == BLACK) {
            //если ребенок красный перекрашиваем в черный
            if (child->color == RED) {
                child->color = BLACK;
            }
            //если ребенок черный вызываем балансировку для ребенка
            else {
                delete_case1(tr, child);
            }
        }
        //если удаляемый узел красный удаляем

        delete n;
    }
    //случай 3: удаляемый узел не имеет детей
    else {
        //если узел черный вызываем балансировку
        if (n->color == BLACK) {
            delete_case1(tr, n);
        }

        //удаляем узел
        if (n->parent) {
            if (n == n->parent->left)
                n->parent->left = NULL;
            else
                n->parent->right = NULL;
        }
        else {
            tr = NULL;  //удаляем корень
        }

        delete n;
    }
}

\end{Verbatim}



\newpage

\section{Анализ}

Красно-чёрное дерево --- это самобалансирующееся двоичное дерево поиска, которое гарантирует логарифмическую высоту при любом порядке вставки и удаления элементов. Из свойств красно-чёрного дерева следует, что его высота \(h\) с \(n\) узлами удовлетворяет неравенству:
\[
h \leq 2\log_2(n + 1) = O(\log n)
\]

Пусть \(n\) — количество узлов в дереве. Все рассматриваемые операции имеют временную сложность, зависящую от высоты дерева, которая всегда составляет \(O(\log n)\).

\subsection*{Вставка элемента (insert)}

Функция \texttt{insert} выполняет вставку нового элемента и состоит из двух этапов.

\textbf{Этап 1. Поиск места вставки:}
\begin{verbatim}
while (y) {
    if (x < y->inf) y = y->left;
    else if (x > y->inf) y = y->right;
    else return;
}
\end{verbatim}
Итеративный спуск от корня до места вставки. Количество итераций не превышает высоты дерева \(h = O(\log n)\).

\textbf{Этап 2. Балансировка (insert\_case1 – insert\_case5):}
Функции балансировки восстанавливают свойства красно-чёрного дерева. Алгоритм поднимается вверх по дереву, пока родитель узла красный. На каждой итерации возможны случаи:
\begin{itemize}
    \item Дядя красный: перекрашивание родителя, дяди и деда. Указатель поднимается на два уровня вверх.
    \item Дядя чёрный: выполняется один или два поворота (каждый за \(O(1)\)) и перекрашивание, после чего цикл завершается.
\end{itemize}
В худшем случае выполняется \(O(h)\) итераций подъёма. Каждая итерация требует \(O(1)\) операций (сравнения, перекрашивания, повороты). Поворотов выполняется не более двух за всю балансировку.

\textbf{Общая временная сложность вставки:}
\[
T_{\text{insert}}(n) = O(\log n) + O(\log n) = O(\log n)
\]

\subsection*{Поиск элемента (find)}

Функция \texttt{find} выполняет стандартный поиск в двоичном дереве:
\begin{verbatim}
if (!tr || x == tr->inf) return tr;
if (x < tr->inf) return find(tr->left, x);
else return find(tr->right, x);
\end{verbatim}
На каждом шаге рекурсии происходит спуск на один уровень вниз. В худшем случае искомый элемент находится в листе или отсутствует — количество рекурсивных вызовов равно высоте дерева \(h = O(\log n)\).

\[
T_{\text{find}}(n) = O(\log n)
\]

\subsection*{Поиск минимума и максимума (find\_min, find\_max)}

Функции реализованы итеративно:
\begin{verbatim}
while (node->left) node = node->left;  // для минимума
while (node->right) node = node->right; // для максимума
\end{verbatim}
Количество итераций не превышает высоты дерева \(h = O(\log n)\).

\[
T_{\min/\max}(n) = O(\log n)
\]

\subsection*{Удаление элемента (delete\_node)}

Функция \texttt{delete\_node} выполняет удаление узла и последующую балансировку.

\textbf{Этап 1. Поиск удаляемого узла:}
Цикл поиска в худшем случае требует \(O(\log n)\) времени.

\textbf{Этап 2. Обработка случаев:}
\begin{itemize}
    \item \textbf{Два потомка:} Находится преемник (максимум в левом поддереве или минимум в правом) за \(O(\log n)\). Значения копируются, после чего удаляется узел-преемник (имеющий не более одного потомка).
    \item \textbf{Один потомок:} Замена узла ребёнком (\texttt{replace\_node}) — \(O(1)\). Если удаляемый узел чёрный, требуется балансировка.
    \item \textbf{Нет потомков (лист):} Если узел чёрный, требуется балансировка. Удаление ссылки у родителя — \(O(1)\).
\end{itemize}

\textbf{Этап 3. Балансировка после удаления (delete\_case1 – delete\_case6):}
Балансировка выполняется только при удалении чёрного узла. Алгоритм проходит вверх по дереву, восстанавливая чёрную высоту. В худшем случае путь от удаляемого узла до корня составляет \(O(\log n)\). На каждом шаге возможно выполнение поворотов (не более трёх за всю балансировку) и перекрашиваний — \(O(1)\).

\textbf{Общая временная сложность удаления:}
\[
T_{\text{delete}}(n) = O(\log n) + O(\log n) + O(\log n) = O(\log n)
\]

\subsection*{Пространственная сложность (расход памяти)}

\textbf{Память для хранения дерева:}
Каждый узел структуры \texttt{tree} содержит:
\begin{itemize}
    \item \texttt{int inf} — 4 байта
    \item \texttt{bool color} — 1 байт (с учётом выравнивания до 4 байт)
    \item \texttt{tree* left} — указатель (4 или 8 байт)
    \item \texttt{tree* right} — указатель (4 или 8 байт)
    \item \texttt{tree* parent} — указатель (4 или 8 байт)
\end{itemize}
Итого на один узел: \(4 + 4 + 3 \times \text{sizeof(pointer)} = 8 + 3p\) байт (на 64-битной системе — 32 байта с учётом выравнивания). Для \(n\) узлов:
\[
M_{\text{storage}}(n) = O(n)
\]

\textbf{Дополнительная память на стеке вызовов:}
Рекурсивная функция \texttt{find} использует стек глубиной до \(h = O(\log n)\). Остальные функции (\texttt{insert}, \texttt{delete\_node}, \texttt{find\_min}, \texttt{find\_max}) реализованы итеративно и не используют дополнительную память на стеке вызовов.

\[
M_{\text{stack}}(n) = O(\log n)
\]

\textbf{Общая пространственная сложность:}
\[
M(n) = O(n) + O(\log n) = O(n)
\]

\end{document}
